\subsection{Matrix-Matrix Multiplication}
\hrulefill

\paragraph*{Definition}
$\forall A \in \mathbb{R}^{m \times n}, B \in \mathbb{R}^{n\times p} \implies AB \in \mathbb{R}^{m \times p} = \begin{bmatrix}
    A\vec{b_1} & A\vec{b_2} & \ldots & A\vec{b_p} \\
\end{bmatrix}$ where $\vec{b_i}$ is the $i$-th column of $B$.

\paragraph*{Properties}
\begin{itemize}
    \item $AB \neq BA$
    \item $A(BC) = (AB)C$
    \item $A(B+C) = AB + AC$
    \item $(A+B)C = AC + BC$
    \item $k(AB) = (kA)B = A(kB) \text{ for any scalar } k$
    \item $(AB)^T = B^T A^T$
\end{itemize}

\subsubsection*{Block Multiplication}
\paragraph*{Definition}
We can partition $A$ and $B$ into blocks and perform multiplication on the blocks. For example, if we partition $A$ into four blocks and $B$ into four blocks, we can write:
\[A = \begin{bmatrix}
    A_{11} & A_{12} \\
    A_{21} & A_{22}
\end{bmatrix}, \quad B = \begin{bmatrix}
    B_{11} & B_{12} \\
    B_{21} & B_{22}
\end{bmatrix}\]
Then, the product $AB$ can be computed as:
\[AB = \begin{bmatrix}
    A_{11}B_{11} + A_{12}B_{21} & A_{11}B_{12} + A_{12}B_{22} \\
    A_{21}B_{11} + A_{22}B_{21} & A_{21}B_{12} + A_{22}B_{22}
\end{bmatrix}\]
This method is particularly useful for large matrices, as it allows us to break down the multiplication into smaller, more manageable pieces.

\paragraph*{Theorem}
If $A, B \in \mathbb{R}^{n \times n}$ are partitioned into comparable blocks, then the product $AB$ can be computed by matrix multiplication
using blocks as entries. In other words, if $A$ and $B$ are partitioned into blocks of size $m \times m$, then the product $AB$ can be computed by multiplying the corresponding blocks and summing the results.

\paragraph*{Multiplying Diagonal Matricies}
If $A$ and $B$ are diagonal matrices, then their product is also a diagonal matrix. Specifically, if $A = \text{diag}(a_1, a_2, \ldots, a_n)$ and $B = \text{diag}(b_1, b_2, \ldots, b_n)$, then their product is given by:
\[AB = \begin{bmatrix}
    a_1 b_1 & 0 & \ldots & 0 \\
    0 & a_2 b_2 & \ldots & 0 \\
    \vdots & \vdots & \ddots & \vdots \\
    0 & 0 & \ldots & a_n b_n
\end{bmatrix}\]